\section{Final remarks}
\label{sec:final-remarks}

The traditional workflow may yield high error rates for entity resolution based on more or
less free texts, because hand-crafted preprocessing rules and traditional scoring models
cannot handle natural language as effectively as structured linkage key values.

The transformer-only workflow, on the other hand, may not be scalable in many
applications, because the cross-encoder model needs to handle a potentially large number
of text pairs directly. Assessment of the associated linkage errors presents another
challenge in practice, particularly when the available labelled sample is used both for
fine-tuning and for error assessment.

The proposed workflow replaces hand-crafted text processing by automated bi-encoder
embedding and refines the steps of scoring and classification (Figure
\ref{fig:combined-workflow}). This allows one to deploy the cross-encoder model to only
a subset of selected cases, while capitalising on the cost-effective baseline performance
that is achievable by the traditional generative or discriminative models. The combined use of the models provides a practical mechanism for controlling
the tradeoff between linkage accuracy and computational cost.

The empirical results illustrate the different roles of the components in this workflow.
Bi-encoder embeddings can improve on traditional string comparisons, especially for
messier texts such as those in Abt-Buy. Fine-tuning can further improve the bi-encoder
representation, although the effect depends on both the dataset and the scoring model.
The fine-tuned cross-encoder yields a favourable balance between FLR and MMR on both
datasets, but at a higher computational cost. The discrimination-gap results show that cross-encoder evaluations can be
concentrated on the harder cases, while many of the easier cases can remain with the
traditional scoring-classification model. 
\rev{The explicit combined-procedure experiment
makes this tradeoff concrete: using the 50th percentile of the training-record gaps as the
routing cutoff reduces cross-encoder pair evaluations by 51.5\% for DBLP-Scholar and
41.0\% for Abt-Buy, with pooled FLR/MMR of 0.017/0.024 and 0.042/0.145,
respectively.}

\rev{The empirical findings are based on two benchmark datasets, and the extent to which the observed patterns generalise to other entity-matching settings remains to be assessed.}

The issues of linkage error assessment have also been studied. Test-set error estimates
refer directly to the models fitted on the corresponding training samples, and need not
equal the errors of models subsequently retrained on a larger labelled sample. The
subsample experiments show that the discrepancy is not of a common direction for both
FLR and MMR. A workable approach is therefore to assess the combined procedure under
the same routing rule that will be used for classification, using the appropriate
denominators of FLR and MMR for combining the component errors. If the component
models are subsequently retrained on the complete labelled sample, additional assumptions
or evidence are needed to relate these test-set estimates to the errors of the retrained
models.

Let us close by outlining three topics for future research on using transformer models for
entity resolution.

The first topic concerns embedding. As noted before, a weakness of the bi-encoder is
that it must encode each text $\tau_i$ without ``seeing'' the other similar texts, including
the likely matched item (Table \ref{tab:silverman-example}). \citet{morris2025contextual} propose a contextual document embedding (CDE) model,
which first uses a sample of texts from the same domain to create a domain-specific
contextual sequence, say $\xi$, and then encodes each concatenated item
$(\xi,\tau_i)$, instead of just $\tau_i$ on its own. The embedding
$\mathbf{x}(\xi,\tau_i)$ may improve on $\mathbf{x}(\tau_i)$ given an appropriate
$\xi$, while using a common $\xi$ for every item keeps the computation scalable.
Whether the particular CDE model, or other ideas of contextual embedding, can be
successfully implemented for entity resolution is clearly worth studying.

Secondly, scoring-classification in the unsupervised setting is a topic we did not consider
in this paper. In particular, fine-tuning of the cross-encoder model would remain
necessary if it is to be adapted to the texts of a specific application. One possibility is
to create a pseudo match sample $\widetilde{s}_M$, where a traditional model is used to
obtain automatically a link set for $A_s$ with a sufficiently low estimated FLR, and use
it as if it were a labelled match sample $s_M$. How reliable such pseudo-labels must be,
and how their errors propagate through fine-tuning and subsequent classification, would
need to be investigated. More generally, the application of cross-encoder models in the
unsupervised setting remains an open question.

Finally, the workflow configuration itself may be refined. In particular, one can study how
best to create separate case streams for link classification, including the possibilities
discussed in Section \ref{sec:combining-models}. The choice is intrinsically connected to
the accuracy-cost tradeoff. Computational cost can be evaluated relatively directly, but
the associated linkage error rates are more difficult to assess, both conceptually and
practically. We have proposed a workable approach in the supervised setting and a setup
for exploring the potential discrepancy between subsample and larger-sample error
assessment in Section \ref{sec:subsample-bias}. Corresponding methods for error
assessment in the unsupervised setting are still needed when entity resolution involves
the use of transformer models.